\documentclass[12pt]{article}
\usepackage{amsmath, amssymb}
\usepackage[margin=1in]{geometry}
\setlength{\parskip}{6pt}
\title{QUBO Formulation: Recommender-System Feature Selection Problem}
\date{}
\begin{document}
\maketitle

\subsection*{Recommender-System Feature Selection Problem}

\paragraph{Problem Statement.}
Given a set of $N$ candidate features for a recommender model, each feature $i$ is associated with a utility score $u_i$, and each pair of features $(i,j)$ is associated with a redundancy or complexity penalty $r_{ij}$. The task is to select exactly $K$ features such that the total utility is maximized while the pairwise redundancy penalties are minimized.

\paragraph{Decision Variables.}
For each candidate feature $i \in \{0, 1, \dots, N-1\}$, introduce a binary decision variable $x_i \in \{0, 1\}$, where $x_i = 1$ indicates that feature $i$ is selected, and $x_i = 0$ indicates that it is not selected.

\paragraph{QUBO Derivation.}
\textbf{Step 1.} The original objective is to maximize $\sum_{i=0}^{N-1} u_i x_i - \sum_{0 \le i < j \le N-1} r_{ij} x_i x_j$. Converting to a minimization problem for the QUBO framework requires negating the objective:
\begin{align}
\min \left( -\sum_{i=0}^{N-1} u_i x_i + \sum_{0 \le i < j \le N-1} r_{ij} x_i x_j \right).
\end{align}

\textbf{Step 2.} The cardinality constraint requires that exactly $K$ features are selected:
\begin{align}
\sum_{i=0}^{N-1} x_i = K.
\end{align}

\textbf{Step 3.} To enforce the constraint, we introduce a quadratic penalty term with weight $\lambda > 0$:
\begin{align}
P(\mathbf{x}) &= \lambda \left( \sum_{i=0}^{N-1} x_i - K \right)^2 \\
&= \lambda \left( \sum_{i=0}^{N-1} x_i^2 - 2K \sum_{i=0}^{N-1} x_i + K^2 \right).
\end{align}
Using the idempotent property of binary variables, $x_i^2 = x_i$, the expression simplifies to:
\begin{align}
P(\mathbf{x}) &= \lambda \left( \sum_{i=0}^{N-1} x_i - 2K \sum_{i=0}^{N-1} x_i + K^2 \right) \\
&= \sum_{i=0}^{N-1} \lambda(1-2K)x_i + \sum_{i \neq j} \lambda x_i x_j + \lambda K^2.
\end{align}
Since $x_i x_j = x_j x_i$, the double sum over distinct indices can be rewritten as twice the sum over ordered pairs $i < j$:
\begin{align}
P(\mathbf{x}) = \sum_{i=0}^{N-1} \lambda(1-2K)x_i + 2 \sum_{0 \le i < j \le N-1} \lambda x_i x_j + \lambda K^2.
\end{align}

\textbf{Step 4.} Combining the negated objective and the penalty term yields the total QUBO Hamiltonian $H(\mathbf{x})$:
\begin{align}
H(\mathbf{x}) &= \sum_{i=0}^{N-1} \left( \lambda(1-2K) - u_i \right) x_i + \sum_{0 \le i < j \le N-1} \left( r_{ij} + 2\lambda \right) x_i x_j + \lambda K^2.
\end{align}

\textbf{Step 5.} Reading off the coefficients from $H(\mathbf{x})$ gives the symmetric QUBO matrix entries and constant offset:
\begin{align}
Q_{i,i} &= \lambda(1-2K) - u_i, \quad \forall i \in \{0, \dots, N-1\}, \\
Q_{i,j} &= r_{ij} + 2\lambda, \quad \forall 0 \le i < j \le N-1, \\
Q_{j,i} &= Q_{i,j}, \quad \forall 0 \le i < j \le N-1, \\
c &= \lambda K^2.
\end{align}

\paragraph{Penalty Weight Justification.}
The penalty weight must be chosen sufficiently large to ensure that any violation of the cardinality constraint is strictly more costly than the maximum possible improvement in the original objective. Specifically, we require:
\begin{align}
\lambda > \max_{i} |u_i| + \sum_{0 \le i < j \le N-1} |r_{ij}|.
\end{align}
This bound guarantees that the penalty incurred by deviating from exactly $K$ selected features dominates any potential gain from the utility or redundancy terms.

\paragraph{Correctness Argument.}
Minimizing $H(\mathbf{x})$ over $\{0,1\}^N$ recovers the optimal feature selection because (i) all feasible solutions satisfying $\sum_i x_i = K$ incur zero penalty and thus evaluate exactly to the negated original objective, and (ii) any infeasible solution incurs a penalty strictly greater than the maximum possible objective value. Consequently, the global minimum of $H(\mathbf{x})$ must be feasible, and among feasible solutions, it must minimize the negated objective, which is equivalent to maximizing the original utility-minus-redundancy score.

\paragraph{Illustrative Example.}
Consider a concrete instance with $N=3$ candidate features and $K=2$ required selections. Let the utility scores be $u_0=10$, $u_1=8$, $u_2=5$, and the pairwise redundancy penalties be $r_{01}=2$, $r_{02}=3$, $r_{12}=1$. We choose $\lambda=20$, which satisfies the justification bound. Substituting these values into the general formulas yields the following concrete terms:
\begin{align}
\text{Linear terms: } & (20(1-4)-10)x_0 + (20(1-4)-8)x_1 + (20(1-4)-5)x_2 = -70x_0 - 68x_1 - 65x_2, \\
\text{Quadratic terms: } & (2+40)x_0x_1 + (3+40)x_0x_2 + (1+40)x_1x_2 = 42x_0x_1 + 43x_0x_2 + 41x_1x_2, \\
\text{Offset: } & 20 \cdot 2^2 = 80.
\end{align}
The resulting Hamiltonian is:
\begin{align}
H(\mathbf{x}) = -70x_0 - 68x_1 - 65x_2 + 42x_0x_1 + 43x_0x_2 + 41x_1x_2 + 80.
\end{align}
Evaluating $H(\mathbf{x})$ over the three feasible binary vectors (where $x_0+x_1+x_2=2$):
\begin{align}
H(1,1,0) &= -70 - 68 + 42 + 80 = -16, \\
H(1,0,1) &= -70 - 65 + 43 + 80 = -12, \\
H(0,1,1) &= -68 - 65 + 41 + 80 = -12.
\end{align}
The global minimum is $-16$ at $\mathbf{x}^* = (1,1,0)$, corresponding to selecting features $0$ and $1$. The maximum original score is $10 + 8 - 2 = 16$, confirming the formulation.

\end{document}